\section{System Characteristics \& Limitations}
\label{sec:results}

\subsection{Empirical Characteristics}

\begin{table}[t]
	\centering
	\begin{tabular}{lcc}
		\hline
		\textbf{Design} & \textbf{Relative token usage} & \textbf{Relative wall-clock runtime} \\
		\hline
		VerCore (Design Conductor 1.0) & 1.0x & 1.0x \\
		Optimized FP32 add/sub & 2.1x & 0.75x \\
		In-switch allreduce & 0.8x & 0.3x \\
		AES encryptor & 4x & 4x \\
		VerTQ & 12x & 6.7x \\
		\hline
	\end{tabular}
	\caption{Relative token usage and wall-clock runtime by design.}
	\label{tab:relative-usage-runtime}
\end{table}

\begin{figure}
	\centering
	\begin{tikzpicture}
	\pie[
	text=legend,
	radius=3,
	color={blue!40, orange!60, green!50, red!45}
	]{
		6.3/Architecture,
		28.2/Low-level modules,
		32.6/High-level modules,
		32.9/Final integration
	}
	\end{tikzpicture}
	\caption{Token usage breakdown (dollar normalized)}
	\label{fig:tok_breakdown}
\end{figure}

See table \ref{tab:relative-usage-runtime} for relative token usage statistics (these are normalized to dollar token cost as comparing raw input, output, cached input, etc. is relatively meaningless).

Figure \ref{fig:tok_breakdown} shows the token usage breakdown. Not unlike a human team, verification and achieving timing closure are the most expensive components.

\subsection{Limitations}

\subsubsection{Overly Methodical}

We find that the system can be overly methodical or over-cautious when dealing with tasks that are actually as simple as they seem. This manifests when asking the system simple queries -- for example, how many floating-point units there are in the design. We observe cases where instead of responding quickly with the obvious answer it already knows, the system will do an extensive codebase review, checking and re-checking its work many times, only to come to the same answer after keeping the user waiting for 20 minutes.

We are unsure if this is due to the new frontier models the updated harness uses or if it's an inherent harness artifact. 

\subsubsection{Backend vs. uArch fixes}

We observed cases where Design Conductor would propose overly complex RTL changes to problems that could be resolved with simple backend script fixes.

\subsubsection{Goal-setting}

We found that in terms of overall constraint and goal-setting, Design Conductor was sometimes too aggressive -- for example, aiming for too high a clock frequency target. 

\subsubsection{Human Review Bottleneck}

While Design Conductor 2.0 is capable of handling very large end-to-end tasks, we find, as many teams on the agent frontier have, that human review is the major bottleneck. The only ``solution'' we have found to this is for the human designers to think through the key design constraints (PPA) upfront.